Sim. Em 64 dos 75 municípios, a proporção de domicílios com acesso à internet (D0330 igual a 1) supera a proporção com esgotamento considerado adequado pelo critério do exercício, que classifica como adequados os valores da variável D0250 até 3. Em Sergipe, as proporções são de 81,34\% para internet e 67,04\% para esgotamento adequado, diferença de 14,30 pontos percentuais. A Tabela~\ref{tab:item2-7-internet-esgoto} mostra os municípios com as maiores diferenças positivas e aqueles em que as proporções são mais próximas, mas ainda há mais cobertura de internet. Há também diferenças no modo de provisão: a família pode contratar internet de um provedor quando existe cobertura; já ampliar a rede coletora de esgoto até as residências exige infraestrutura coletiva e decisões de planejamento e investimento público. Além disso, o critério adotado considera adequada a fossa séptica ou fossa-filtro não ligada à rede (valor 3 de D0250), de modo que a medida não representa apenas o alcance da rede coletora. O contraste evidencia desigualdade territorial entre serviços essenciais, mas os dados agregados por município não permitem concluir que as famílias tenham escolhido internet em vez de saneamento.
